\chapter{Signals Lean library}

This blueprint indexes selected definitions and results in the Signals Lean
library. The library is organized into modules covering mathematical and
physical models, signal propagation, measurement, and applications.

Only Lean declarations marked with LeanArchitect's \texttt{@[blueprint]}
attribute are emitted as formal nodes. The extracted module index is generated
from the Lean sources before each blueprint build.

Declaration docstrings document the Lean API. Blueprint statements, titles and
proof summaries are supplied explicitly in the attributes; ordinary declaration
docstrings do not populate the blueprint statement. Statement and proof
dependencies are inferred from the corresponding Lean terms.

Readiness markers report sorry-free formal declarations, not experimental
validation, calibrated data or the realization of a physical model.

\input{../../.lake/build/blueprint/library/Signals.tex}
\input{../../.lake/build/blueprint/library/SignalsTests.tex}
\input{../../.lake/build/blueprint/library/SignalsPending.tex}
\input{../../.lake/build/blueprint/library/SignalsPendingTests.tex}

\section{Finite routing and weighted measurements}

The rotation-system and weighted-network layers are distinct. Neither supplies
a planar disk embedding, reducedness or a general disjoint-path positivity
theorem. Fixed-point decorations and signed-pair compatibility remain explicit.

\inputleanmodule{Signals.Plabic}

\section{Ordered minors and geometric representatives}

These nodes describe finite real matrix representatives, not a Grassmannian
quotient. Positive-condition records need an available ordered selection for
rank promotion; nonnegative full-rank records retain a nonzero-minor witness.
Two-step identities do not identify path weights with complex optical amplitudes.

\inputleanmodule{Signals.Geometry}

\section{Scalar propagation and amplitude controls}

Finite sums and planar integrals use the chosen scalar approximation.
Regularity and integrability guards exclude totalized singular inputs;
intensity remains uncalibrated. Comparison records do not prove independently
evaluated models agree or solve an exact diffraction boundary problem.

\inputleanmodule{Signals.Huygens}

\section{Measure-derived quadrature and refinement}

Error rates require the stated measure, integrability, Lipschitz and mesh
hypotheses. The nodes do not construct a rectangular mesh or establish
kernel-specific regularity, a Fraunhofer limit or measured convergence.

\inputleanmodule{Signals.Quadrature}

\section{Affine forms and oriented residues}

These are fixed-chart interval, triangle and square examples. Residues are
punctured limits; diagonal continuation is separate from off-diagonal
triangulation equality. General projective canonical forms, amplituhedron
forms and an optical evaluator remain outside this layer.

\inputleanmodule{Signals.CanonicalForms}

\section{Existing model and propagation nodes}

\inputleanmodule{SignalsPending}
\inputleanmodule{Signals.Propagation}